\documentclass[11pt]{article}
\usepackage[margin=1in]{geometry}
\usepackage{graphicx}
\usepackage{amsmath}
\usepackage{booktabs}
\usepackage{siunitx}
\usepackage{xcolor}
\usepackage{caption}
\usepackage[colorlinks=true,linkcolor=blue!55!black,urlcolor=blue!55!black]{hyperref}
\captionsetup{font=small,labelfont=bf}

% Figures live one level up, in figures/ and event_displays/
\graphicspath{{../figures/}{../event_displays/}}

% Robust include: compiles even if a figure file is renamed/missing.
\newcommand{\pkgfig}[2][width=0.9\textwidth]{%
  \IfFileExists{../figures/#2}{\includegraphics[#1]{#2}}{%
  \IfFileExists{../event_displays/#2}{\includegraphics[#1]{#2}}{%
  \fbox{\parbox{0.8\textwidth}{\centering\vspace{2em}\texttt{MISSING FIGURE: \detokenize{#2}}\vspace{2em}}}}}}

% Key-number box (kept to base packages: no tcolorbox installed).
\definecolor{keybg}{RGB}{235,242,250}
\definecolor{keyframe}{RGB}{70,110,160}
\newsavebox{\keyboxsave}
\newenvironment{keybox}[1]{%
  \begin{lrbox}{\keyboxsave}\begin{minipage}{0.88\textwidth}%
  {\small\bfseries\textcolor{keyframe}{#1}}\par\smallskip
}{%
  \end{minipage}\end{lrbox}%
  \begin{center}\setlength{\fboxsep}{8pt}%
  \fcolorbox{keyframe}{keybg}{\usebox{\keyboxsave}}\end{center}}

\newcommand{\um}{\,\si{\micro\meter}}
\newcommand{\detA}{Detector~A (mx17\_3)}
\newcommand{\detB}{Detector~B (mx17\_2)}

\title{\vspace{-1.5em}The MX17 Micromegas detectors under cosmic muons\\
\large Performance summary of the June 2026 Saclay cosmic-bench campaign\\
\normalsize Reference material for the detector-construction conference talk}
\author{MX17 cosmic-bench analysis}
\date{July 2026 --- preliminary}

\begin{document}
\maketitle
\vspace{-1.5em}

\begin{quote}\small\itshape
2026-07-14 note: the M3 reference recipe tightened to $\chi^2<1.0$,
$N_\mathrm{clus}=4$ (was $\chi^2<5$, $N_\mathrm{clus}\ge3$) after finding the
position residual was reference-limited, not chamber-limited. All numbers in
this document are re-measured on the new recipe (efficiency, position
resolution, and angular resolution all updated fleet-wide); the gas/drift-velocity
and edge/fringe-field sections predate the change and are not expected to shift
materially. See \texttt{M3\_CUT\_AND\_ACTIVE\_AREA\_NOTE.md} in the analysis
repository for the full investigation.
\end{quote}

\begin{abstract}
\noindent This document summarizes, in plain language, how the five MX17
resistive-strip Micromegas detectors behave when they detect cosmic muons, based
on the June~2026 test-bench campaign at Saclay. It is written as source material
for a conference talk on the construction of these detectors: every headline
number is quotable, every figure exists as a standalone PNG and PDF in the
accompanying \texttt{figures/} and \texttt{event\_displays/} directories, and a
starter slide deck is provided. In one sentence: the detectors work very well ---
the best chamber detects \(\sim\)96\% of the muons that cross its central region,
locates each crossing to about $0.47\,$mm (better than the $0.78\,$mm strip
pitch), measures the track direction to about $1.7^\circ$, times the event to
about 33\,ns ($\approx$1\,mm of electron drift), and its built-in resistive
spark protection works exactly as designed: discharges are rare, localized,
self-quenching, and cause \emph{no measurable dead time}.
\end{abstract}

\tableofcontents
\newpage

%=====================================================================
\section{What this document is, and the one-slide summary}

Five MX17 Micromegas chambers (labelled A--E at the experiment) were
characterized with cosmic muons on the Saclay test bench during June~2026.
The analysis behind this summary is mature: every detector was measured with
the same reference telescope, the same reconstruction, and the same definitions,
so the numbers below can be compared detector-to-detector and quoted directly.

\begin{keybox}{Headline numbers (best detector, A, at its operating point)}
\begin{tabular}{@{}ll@{}}
Muon detection efficiency & $\mathbf{92.9\pm0.2\,\%}$ over the active area (operating; $96.6\,\%$ spark-free); $\sim$\textbf{96\,\%} in the core \\
Position resolution & $\mathbf{\sigma \approx 0.47\,mm}$ (sub-pitch; strip pitch 0.78\,mm)\\
Track-angle resolution & $\mathbf{\approx 1.7^\circ}$ per plane, at all track angles \\
Event-time resolution & $\mathbf{33\,ns}$ ($\approx$1\,mm of drift); 37.7\,ns absolute vs.\ trigger \\
Discharge (spark) behaviour & rare (3.9\,\% of crossings), self-quenching, \textbf{zero dead time} \\
\end{tabular}
\end{keybox}

\noindent The five detectors are not identical in performance --- two are
excellent, two were operated past their sparking optimum during the tests, and
one has a gas-gain problem --- and Section~\ref{sec:fleet} gives the honest
per-detector picture with the reasons. All of the differences are consistent
with operating-point and gas effects, not with construction defects.

\begin{figure}[htbp]
  \centering
  \pkgfig[width=0.6\textwidth]{01-mx17-layout-topdown.png}
  \caption{Where the chambers are headed: at the experiment, four of the MX17
  Micromegas surround the central target in a pinwheel layout. The bench
  campaign described here characterized each chamber individually with cosmic
  muons before installation.}
  \label{fig:pinwheel}
\end{figure}

\paragraph{How to use this package.} The directory layout is:
\texttt{figures/} (every plot as PNG + PDF, with \texttt{FIGURE\_GUIDE.md}
describing each one), \texttt{event\_displays/} (fresh, presentation-quality
single-muon event pictures), \texttt{slides/} (a starter PowerPoint deck),
\texttt{source\_reports/} (the detailed technical write-ups behind each topic),
and this report.

%=====================================================================
\section{The detectors and how they were tested}
\label{sec:bench}

\subsection{The device under test, as operated}

Each MX17 chamber is a resistive-strip Micromegas with a two-dimensional strip
readout: 512~strips in $X$ and 512~strips in $Y$ at a pitch of
$0.78\,$mm, giving a $\sim$40$\times$40\,cm$^2$ active area. A pixelated top
layer divides the avalanche charge between the two strip layers, so a single
avalanche is seen by \emph{both} coordinates at once (Section~\ref{sec:balance}
shows how evenly). Above the amplification mesh sits a \textbf{30\,mm
conversion-and-drift gap}: a cosmic muon crossing the chamber liberates
ionization electrons along its path through this gas volume, and these drift
down to the mesh where they are amplified and read out.

Operating configuration on the bench:
\begin{itemize}
  \item \textbf{Gas:} Ar/isobutane 95/5 at atmospheric pressure, flushed
        continuously (the campaign also quantified what happens when moisture
        and air contaminate the gas --- Section~\ref{sec:gas}).
  \item \textbf{High voltage:} amplification (``resist'') voltage typically
        480--490\,V; drift voltage typically 1000\,V across the 30\,mm gap
        ($\approx$330\,V/cm drift field).
  \item \textbf{Readout:} DREAM front-end electronics, 512~channels per plane,
        sampling each strip at 60\,ns intervals over a 1.92\,$\mu$s window ---
        so every strip records a full pulse \emph{waveform}, with both an
        amplitude and a time.
\end{itemize}

\subsection{The cosmic bench}

\begin{figure}[htbp]
  \centering
  \pkgfig[width=0.62\textwidth]{00-cosmic-bench-schematic.png}
  \caption{The Saclay cosmic bench (side view, heights to scale): trigger
  scintillators above and below, two pairs of M3 reference Micromegas, and the
  two MX17 test slots in between.}
  \label{fig:bench}
\end{figure}

The bench provides two independent references against which the chambers are
judged:
\begin{itemize}
  \item \textbf{A reference tracker (``M3'')}: a telescope of multiplexed
        Micromegas planes above and below the test slots. For every triggered
        cosmic muon it reconstructs the muon's straight line in space,
        independently of the detector under test. Every performance number in
        this document is defined against these reference tracks: the track says
        where the muon truly went; the chamber under test either saw it there
        or it did not.
  \item \textbf{A scintillator trigger}: a top+bottom scintillator-paddle
        coincidence ($\sim$5\,ns precision) starts the readout, providing the
        absolute time reference used for the timing measurement.
\end{itemize}

Alignment of each chamber to the telescope converged to better than a
millimetre in every case, and the recovered geometry (rotation
$\approx$89.5$^\circ$, heights consistent with the mechanical slot positions)
matches the as-built bench.

\subsection{The definitions used everywhere}
\label{sec:defs}

\begin{description}
  \item[Efficiency] For every clean reference muon track, project it to the
  chamber plane and ask: did the chamber reconstruct a hit (with both $X$ and
  $Y$ coordinates) within 5\,mm of that point? The quoted efficiency is the
  fraction of ``yes''. A muon for which the chamber recorded nothing at all
  still counts in the denominator --- there is no hidden normalization.
  \item[Position resolution] The width (Gaussian core, $\sigma$) of the
  difference between the chamber's reconstructed position and the reference
  track's position. Note this still \emph{includes} the small pointing error of
  the reference telescope itself, so all position resolutions quoted here are
  slightly conservative (upper limits on the true detector resolution).
  \item[Angular resolution] The spread ($\sigma$ or the 68\% half-width,
  $\sigma_{68}$) of the difference between the track angle measured by the
  chamber alone (Section~\ref{sec:tpc}) and the reference track angle.
  \item[Spark fraction] Discharges fire most of the detector at once
  ($>$50~strips); they are counted separately as the percentage of muon
  crossings that coincide with such a discharge.
\end{description}

%=====================================================================
\section{How the chamber sees a muon: micro-TPC operation}
\label{sec:tpc}

The 30\,mm drift gap is not just a conversion volume --- it turns each chamber
into a small time-projection chamber (``micro-TPC''). Electrons freed deeper in
the gap arrive at the mesh \emph{later}: at the nominal drift field the
electrons travel at $34\um$/ns, so the 30\,mm gap maps onto a $\sim$0.9\,$\mu$s
spread of arrival times, comfortably resolved by the 60\,ns sampling. Each
fired strip therefore contributes a point (strip position, arrival time)
$=$ (horizontal position, depth), and a straight-line fit through those points
gives the muon's crossing angle \emph{from a single chamber} --- no external
tracker needed.

\begin{figure}[htbp]
  \centering
  \pkgfig[width=0.85\textwidth]{02-microtpc-principle-schematic.png}
  \caption{The micro-TPC principle: an inclined muon (left) produces a ladder
  of strip signals whose arrival times encode depth (right); the slope of the
  ladder is the track angle.}
  \label{fig:principle}
\end{figure}

\subsection{Event displays}

Figures~\ref{fig:evgallery} and~\ref{fig:evwave} are real single-muon events
from detector~A, drawn directly from the raw data. In the strip-vs-time view
the muon is simply a straight line of charge; in the raw waveform view the
diagonal charge stripe across the channel-vs-time plane is the muon track as
the electronics actually record it. Figure~\ref{fig:ev3d} puts the two strip
planes together into a single 3-D picture: the drift time of each strip signal
becomes a depth, so the charge cloud traces the muon's path through the gas ---
and the independent reference telescope track (green) runs straight through it,
a direct visual check of both the reconstruction and the alignment.

\begin{figure}[htbp]
  \centering
  \pkgfig[width=0.95\textwidth]{event_display_gallery.png}
  \caption{Gallery of cosmic-muon events in detector A. Each panel shows one
  muon crossing: every marker is one strip signal, placed at its strip position
  (horizontal) and electron drift time (vertical, equivalently depth in the
  gas), colored by pulse amplitude. The overlaid line is the track fitted from
  this chamber's own data.}
  \label{fig:evgallery}
\end{figure}

\begin{figure}[htbp]
  \centering
  \pkgfig[width=0.95\textwidth]{event_display_waveforms.png}
  \caption{``Anatomy of an event'': the raw DREAM waveforms for one inclined
  muon. The diagonal stripe in the channel-vs-time plane is the muon; its slope
  is the drift-time ladder used for the angle measurement.}
  \label{fig:evwave}
\end{figure}

\begin{figure}[htbp]
  \centering
  \pkgfig[width=0.98\textwidth]{08-det3A-event-3d-gallery.png}
  \caption{Four real events in 3-D, spanning $\sim$6--20$^\circ$. Each marker is
  one strip signal (circle = X plane, square = Y plane) placed at its strip
  position and at the depth implied by its electron drift time, colored by pulse
  amplitude. The green line is the \emph{independent} M3 telescope track,
  extrapolated into the chamber --- it threads the charge cloud in every event.
  These are among the tightest matches (0.1--0.3\,mm at the mesh); across the
  full sample the chamber-to-reference match has a $0.6$--$0.8$\,mm per-axis core
  (Section~\ref{sec:position}), so a \emph{typical} event sits $\sim$0.8\,mm off
  in this 2-D radial sense --- the display events were chosen near the tight end
  to show the geometry cleanly. Any residual fanning of the raw per-strip ladder
  away from the line with depth is the resistive-strip charge-sharing bias
  (Section~\ref{sec:sharing}), removed before angles are quoted --- it is not a
  position error. Single-event versions (with turntable animations) are in
  \texttt{event\_displays\_3d/}.}
  \label{fig:ev3d}
\end{figure}

\subsection{A built-in subtlety: charge sharing between strips}
\label{sec:sharing}

The resistive strips deliberately spread the avalanche charge across
neighbours: a strip receives $\sim$45--52\% of its neighbour's prompt charge
(and a few percent at two strips), with a $\sim$70\,ns delay, and the spread
reaches $\sim$3$\times$ further along the resistive-strip direction than across
it. This was measured directly from vertical tracks and is the same on both
chambers tested --- it is a \emph{design property}, not a defect. It cuts both
ways:
\begin{itemize}
  \item It distorts the naive time ladder (shared charge arrives with the
  neighbour's timing), which biases simple angle fits; the analysis removes
  this with an ``unsharing'' correction, after which the time-based and
  geometry-based measurements agree.
  \item It is repaid in full for position: the sharing acts as a built-in
  charge interpolator between strips, which is precisely why the position
  resolution beats the strip pitch (Section~\ref{sec:position}).
\end{itemize}

\begin{figure}[htbp]
  \centering
  \pkgfig[width=0.9\textwidth]{03-charge-sharing-schematic.png}\\[4pt]
  \pkgfig[width=0.8\textwidth]{18-det3A-charge-sharing-measured.png}
  \caption{Top: how resistive-strip charge sharing works and why it distorts a
  naive time fit. Bottom: the sharing measured in data from vertical tracks ---
  the first neighbour carries $\sim$45\% (X) / 52\% (Y) of the lead strip's
  charge, with a $\sim$70\,ns delay.}
  \label{fig:sharing}
\end{figure}

%=====================================================================
\section{Headline performance of the best chamber (Detector A)}
\label{sec:det3}

All numbers in this section are from \detA{} at its operating point
(490\,V amplification, 1000\,V drift), using the highest-statistics June runs
(22{,}300 clean reference muons in the headline run at the current recipe ---
fewer than the previous 52{,}000 because the stricter reference-track cut trades
statistics for a cleaner reference, see the note on page~1).

\subsection{Detection efficiency}
\label{sec:eff}

\begin{keybox}{Efficiency --- Detector A}
$\mathbf{92.9\pm0.2\,\%}$ of reference muons are reconstructed within 5\,mm over
the full active area --- the operating number, with the 3.9\,\% in-spark
coincidence folded in (spark-free it is $96.6\pm0.1\,\%$); $\sim$\textbf{96\,\%}
in the core ($>$25\,mm from the frame, pre-2026-07-14 figure, not yet
re-verified). Crucially, the chamber produces a signal
for \textbf{100.0\,\%} of muons and a reconstructed track point for 96.0\,\%:
genuine blindness is effectively $0.0\,\%$. The loss off 92.9\,\% is a discharge
coincidence (3.9\,\%) plus an edge/near-miss \emph{position} tail (3.1\,\%), not
dead area.
\end{keybox}

The efficiency is uniform across the chamber face except for the narrow edge
band: Figure~\ref{fig:effmap} shows the map. Three numbers describe the same
chamber and it is worth keeping them straight: \textbf{92.9\,\%} is the
operating efficiency (everything, including the muons that arrive during a
discharge); \textbf{96.6\,\%} is the same measurement on spark-free crossings;
$\sim$\textbf{96\,\%} is the core, away from the $\sim$25\,mm fringe-field edge
band that is present by construction in any bounded drift volume. The 3.1\,\%
that reconstructs \emph{past} the 5\,mm cut is a near-miss position tail
(competing clusters, edge distortion, sub-veto discharges), almost all inside
10\,mm --- loosening the match recovers still more of it. The chamber
essentially never fails to \emph{see} a muon; the question is only how close
the reconstructed point lands.

\begin{figure}[htbp]
  \centering
  \pkgfig[width=0.98\textwidth]{20-det3A-efficiency-map.png}\\[4pt]
  \pkgfig[width=0.85\textwidth]{21-det3A-efficiency-breakdown.png}
  \caption{Top: Detector A efficiency map --- fraction of reference muons
  matched within 5\,mm vs.\ position (left), fraction with any signal at all
  (middle, uniformly $\sim$100\%), and track statistics (right). Bottom: the
  loss budget --- 92.9\% reconstruct in place; the remainder is near-misses
  (3.1\%) and in-spark coincidences (3.9\%); essentially no muons leave no
  signal.}
  \label{fig:effmap}
\end{figure}

\subsection{Position resolution: better than the strip pitch}
\label{sec:position}

\begin{keybox}{Position resolution --- Detector A}
Gaussian-core residuals against the telescope: $\sigma = 0.44$--$0.47\,$mm
(radial core, $\chi^2<1.0$+$N_\mathrm{clus}=4$ recipe), and the best position
estimator reaches \textbf{0.47/0.54\,mm} ($X$/$Y$) for near-vertical muons ---
\emph{below the 0.78\,mm strip pitch}. Sub-pitch precision is possible only
because the resistive layer shares charge between neighbouring strips,
providing analog interpolation. All values still include the reference
telescope's own pointing error, so the true detector resolution is slightly
better.
\end{keybox}

\begin{figure}[htbp]
  \centering
  \pkgfig[width=0.95\textwidth]{10-det3A-spatial-residuals.png}
  \caption{Detector A position residuals with respect to the reference track
  in $X$ (top row) and $Y$ (bottom row): full distribution, zoom, and Gaussian
  core fit ($\sigma = 0.61/0.73$\,mm on this run, telescope included).}
  \label{fig:residuals}
\end{figure}

\subsection{Track-angle resolution: $\sim$1.7$^\circ$ from a single chamber}
\label{sec:angle}

\begin{keybox}{Angular resolution --- Detector A (hybrid tracking)}
$\sigma_{68} \approx 1.7^\circ$ per readout plane ($1.66^\circ$ in the
near-vertical band, $1.84^\circ$ on the inclined plateau), \emph{at all track
angles} including head-on, with a residual bias of $\approx0.1^\circ$ and
97--98\% of tracks usable. The same reconstruction transferred untouched to
Detector~B gives $2.7^\circ$ (vs.\ $2.3^\circ$ retrained), confirming this is a
property of the detector design, not a tuning artifact -- every chamber's
transfer behaviour matches the same pattern seen before the M3 recipe change.
\end{keybox}

Two methods are combined: for inclined tracks the drift-time ladder
(micro-TPC fit, after the unsharing correction of Section~\ref{sec:sharing});
for near-vertical tracks --- where the time ladder carries no lever arm --- the
shape of the charge deposit itself. The hybrid of the two delivers a uniform
$\sim$1.7$^\circ$ everywhere. In addition, the fitted drift velocity from the
data agrees with the dedicated measurement to 1--3\%, meaning the chambers can
monitor their own gas condition in situ during the experiment.

\begin{figure}[htbp]
  \centering
  \pkgfig[width=0.62\textwidth]{15-det3A-hybrid-angle-correlation.png}\\[4pt]
  \pkgfig[width=0.7\textwidth]{16-det3A-hybrid-angular-resolution-vs-angle.png}
  \caption{Hybrid tracking, Detector A. Top: track angle measured by the
  chamber alone vs.\ the reference telescope --- the events hug the diagonal.
  Bottom: angular resolution vs.\ track angle. The grey curve is the
  \emph{same} method with the head-on step switched off (``high-angle-only'':
  the drift-time fit alone); it tracks the hybrid above $\sim$5$^\circ$ --- they
  are the same estimator there --- and rises to $\approx$5$^\circ$ on
  near-vertical tracks, where a drift-time fit has no lever arm (and there
  covers only $\sim$51\% of them). The hybrid stays uniform at
  $\approx$1.7$^\circ$ across all angles at 97--98\% coverage.}
  \label{fig:angres}
\end{figure}

\subsection{Time resolution: $\approx$1\,mm of drift}
\label{sec:timing}

\begin{keybox}{Time resolution --- Detector A (replicated on B within 10\%)}
Detector time resolution $\sigma_t = 33\,$ns per plane (29\,ns after time-walk
correction), i.e.\ $\approx$1\,mm of electron drift --- measured
telescope-free from the agreement of the two orthogonal strip layers timing
the same drift electrons (inter-plane bias $-1.3$\,ns, consistent with zero).
Absolute event time versus the scintillator trigger: $\sigma_{68} = 37.7\,$ns,
dominated by the detector term itself.
\end{keybox}

The two strip layers sit under one common amplification stage and read the same
avalanche, so their time difference is a pure, assumption-free measure of the
chamber's own timing --- a clean consequence of the two-coordinate design. The
$\approx$1\,mm drift-equivalent timing matches the $\approx$0.6\,mm transverse
resolution: the chamber is not timing-limited relative to its geometry.

\begin{figure}[htbp]
  \centering
  \pkgfig[width=0.55\textwidth]{06-timing-principle-schematic.png}\\[4pt]
  \pkgfig[width=0.9\textwidth]{51-time-resolution-plane-to-plane.png}
  \caption{Top: the built-in timing measurement --- the $X$ and $Y$ layers
  independently timestamp the same drifting electrons. Bottom: their time
  difference gives $\sigma_t = 33$\,ns per plane with zero bias.}
  \label{fig:timeres}
\end{figure}

%=====================================================================
\section{Operating point: high-voltage behaviour}
\label{sec:hv}

Each detector was scanned in amplification voltage. The behaviour is textbook:
efficiency turns on, reaches a plateau, then \emph{falls} at high voltage ---
and the falloff is entirely explained by sparking, not by the detector going
quiet (the fraction of events with any signal stays $\sim$100\% throughout,
while the spark fraction climbs from a few percent to $\sim$50\% exactly where
the efficiency rolls off). Resolution is best at the efficiency optimum.

\begin{keybox}{Operating points found (drift 1000\,V for A/B, 700\,V for C/D)}
\begin{tabular}{@{}lccc@{}}
\toprule
Detector & Peak efficiency & Optimal amplification HV & Plateau behaviour \\
\midrule
A (mx17\_3) & 93.2\,\% & 490--495\,V & broad plateau, $\sigma\approx0.5$--$0.65$\,mm \\
B (mx17\_2) & 94.9\,\% & 485--490\,V & plateau, sparks off above $\sim$510\,V \\
C (mx17\_6) & 78.8\,\% & 480\,V & full turn-on $\to$ plateau $\to$ falloff \\
D (mx17\_7) & 64.1\,\% & 440\,V & full turn-on $\to$ plateau $\to$ falloff \\
\bottomrule
\end{tabular}
\end{keybox}

\begin{figure}[htbp]
  \centering
  \pkgfig[width=0.85\textwidth]{30-det3A-hv-scan-efficiency-sparks.png}
  \caption{Detector A amplification-voltage scan: efficiency (left axis) and
  spark fraction (right axis). The efficiency loss above the plateau is
  spark-induced, not loss of signal.}
  \label{fig:hvscan}
\end{figure}

%=====================================================================
\section{Discharges: the resistive protection works}
\label{sec:sparks}

Because the amplification structure is resistive, a discharge (``spark'') is
expected to stay small, quench itself, and leave the detector immediately
live. The June campaign tested every part of that expectation directly.

\begin{keybox}{Spark behaviour --- measured, Detector A at operating point}
Rate: 3.9\,\% of muon crossings ($0.33$\,Hz, Poisson --- i.e.\ random, no
bursts). Mostly muon-induced (4$\times$ enhancement in coincidence with a
crossing) and concentrated at the chamber edge. Non-propagating: the raw
waveforms show a fast \emph{global baseline step} that recovers within the
1.92\,$\mu$s readout window in 94\% of cases, with the genuine localized charge
confined to $\sim$40--56 strips at one edge. And critically: \textbf{no
measurable dead time} --- reconstruction efficiency and gas gain are flat as a
function of time since the previous spark, so the only loss is the crossing
that coincides with the spark itself (already inside the efficiency numbers).
\end{keybox}

\begin{figure}[htbp]
  \centering
  \pkgfig[width=0.95\textwidth]{43-det3A-spark-waveform-gallery.png}
  \caption{What a discharge actually looks like in the raw data: a fast,
  global, self-quenching baseline excursion (compare the clean muon of
  Fig.~\ref{fig:evwave}). It is over within the 2\,$\mu$s readout window.}
  \label{fig:sparkwf}
\end{figure}

\begin{figure}[htbp]
  \centering
  \pkgfig[width=0.85\textwidth]{44-det3A-spark-deadtime-null.png}
  \caption{The dead-time null measurement: efficiency and pulse height vs.\
  time since the previous discharge are flat --- the chamber is immediately
  live again.}
  \label{fig:deadtime}
\end{figure}

For a construction talk this is arguably the strongest single message: the
resistive spark-protection scheme was verified at the raw-waveform level and
imposes \emph{no recovery-time penalty} on operation.

%=====================================================================
\section{The pixelated layer shares charge evenly between X and Y}
\label{sec:balance}

The top pixelated layer is supposed to route the avalanche charge evenly to the
two strip coordinates. Measured: the fraction of charge on the $X$ layer,
$f = q_X/(q_X+q_Y)$, is centred at 0.491 (A) / 0.534 (B) --- within a few
percent of the ideal 0.5, with the small per-chamber offset attributable to
assembly --- and is \emph{constant} across the chamber face and across track
angles (spread $\sigma_{68}\approx0.07$ per event; true position-to-position
variation only $\sim$2\%). The $X$ and $Y$ charges of a given event are
strongly correlated ($r\approx0.84$--$0.89$): both layers really do see the same
avalanche.

\begin{figure}[htbp]
  \centering
  \pkgfig[width=0.55\textwidth]{05-charge-routing-schematic.png}\\[4pt]
  \pkgfig[width=0.48\textwidth]{60-xy-charge-balance-correlation.png}%
  \hfill
  \pkgfig[width=0.48\textwidth]{61-xy-charge-balance-distribution.png}
  \caption{Top: the chamber stack --- the pixelated resistive top layer routes
  the avalanche charge to both strip layers. Bottom left: total $X$-layer vs.\
  $Y$-layer charge per event ($r=0.89$). Bottom right: the balance fraction
  $f=q_X/(q_X+q_Y)$ is narrow and centred near 0.5 in both chambers measured.}
  \label{fig:balance}
\end{figure}

%=====================================================================
\section{Gas matters: drift velocity, moisture, and attachment}
\label{sec:gas}

The campaign measured the drift velocity at five drift fields and found
$34\pm1.5\um$/ns at the nominal 1000\,V --- and the shape of the
velocity-vs-field curve identifies the gas: it matches simulation (Magboltz)
for Ar/iso 95/5 with $\sim$1\% water. Two practical lessons for operation:

\begin{itemize}
  \item \textbf{Moisture changes the drift speed.} Detector A's gas visibly
  dried from $>$3\% to $\sim$1\% water over the week of running --- directly
  observable in the data. Water content sets the velocity scale.
  \item \textbf{Air (oxygen) eats the signal.} The measured signal decays
  exponentially with drift \emph{distance} (attachment length
  $\lambda\approx15$\,mm at nominal field, consistent with $\sim$1\% air), so
  electrons from the top of the 30\,mm gap partially disappear: the recorded
  charge column is effectively $\sim$23\,mm. This costs amplitude, not
  efficiency, at the June contamination level --- but it is the reason gas
  tightness and flushing matter.
\end{itemize}

\begin{figure}[htbp]
  \centering
  \pkgfig[width=0.85\textwidth]{70-det3A-drift-velocity-vs-magboltz.png}
  \caption{Measured drift velocity vs.\ drift field, compared with Magboltz
  simulation for different gas admixtures: the data identify Ar/iso 95/5 with
  $\sim$1\% water.}
  \label{fig:vdrift}
\end{figure}

\begin{figure}[htbp]
  \centering
  \pkgfig[width=0.9\textwidth]{71-det3A-attachment-amplitude-decay.png}
  \caption{Signal amplitude vs.\ drift \emph{time} (left; curves at different
  drift voltages disagree) and vs.\ drift \emph{distance} (right; all voltages
  collapse onto one exponential, $\lambda\approx18$\,mm): the loss is per
  millimetre drifted --- electron attachment on oxygen from $\sim$1\% air ---
  and it shortens the usable charge column.}
  \label{fig:attach}
\end{figure}

%=====================================================================
\section{The chamber edge}
\label{sec:edge}

The drift field is distorted within $\sim$25--40\,mm of the frame (the fringe
region unavoidable in any bounded drift volume). Consequences, all measured:
the efficiency turns on from 0 to 96\% over the first 25\,mm from the edge;
track angles acquire a $\sim$3$^\circ$ inward tilt in the band (position
measurements remain good --- only angles are affected); and discharges
preferentially seed there. Everything inboard of the band is uniform. For the
experiment this is handled by a simple fiducial cut.

\begin{figure}[htbp]
  \centering
  \pkgfig[width=0.98\textwidth]{80-det3A-edge-fringe-dashboard.png}
  \caption{The full edge study (Detector A). The key panel is the efficiency
  vs.\ distance-to-edge curve (top right): a 25\,mm turn-on band, then a flat
  $\sim$96\% core. Companion panels show the angle tilt and the slowed drift in
  the same band, and that position residuals stay flat to the edge.}
  \label{fig:edge}
\end{figure}

%=====================================================================
\section{The full fleet: five detectors compared}
\label{sec:fleet}

All five chambers were characterized identically. Detector letters are the
experiment labels (A--E); the mx17 numbers are the construction labels.

\begin{keybox}{Fleet summary (June 2026, best runs, identical definitions)}
\begin{tabular}{@{}lrrrrl@{}}
\toprule
Detector & Efficiency & Position $\sigma$ & Angle $\sigma_{68}$ & Spark & Verdict \\
\midrule
A (mx17\_3) & 92.9\,\% & 0.47\,mm & 1.7$^\circ$ & 3.9\,\% & Best performer \\
B (mx17\_2) & 91.3\,\% & 0.46\,mm & 2.3$^\circ$ & 5.0\,\% & Healthy \\
C (mx17\_6) & 57.8\,\% & 0.45\,mm & 3.9$^\circ$\,$^\dagger$ & 23.0\,\% & Spark-limited \\
D (mx17\_7) & 43.1\,\% & 0.59\,mm & 3.4$^\circ$ & 33.4\,\% & Spark-limited \\
E (mx17\_4) & 20.7\,\% & 0.67\,mm & 2.6$^\circ$ & 3.2\,\% & Gain-limited \\
\bottomrule
\end{tabular}\\[4pt]
{\footnotesize Efficiency within 5\,mm of the reference track, full active
area, 2026-07-14 recipe ($\chi^2<1.0$, $N_\mathrm{clus}=4$; was $\chi^2<5$,
$N_\mathrm{clus}\ge3$ -- superseded values: A 88.8\,\%/0.63\,mm/1.75$^\circ$,
B 87.0\,\%/0.64\,mm/2.47$^\circ$, C 55.7\,\%/0.59\,mm/4.14$^\circ$,
D 36.4\,\%/0.86\,mm/3.64$^\circ$, E 20.0\,\%/0.89\,mm/2.70$^\circ$, do not
quote). Position $\sigma$ includes the reference telescope. Angle $\sigma_{68}$
= the hybrid (current-best) tracking, $|\theta|<5^\circ$ band, self-trained,
odd-eid holdout, run \textbf{fleet-wide} (coverage A 98\,\%, B 96\,\%, C 98\,\%,
D 90\,\%, E 79\,\%). The 700\,V spark-limited C/D and gain-limited E are
degraded but measured. $^\dagger$C is low-angle-only: its X-plane develops no
drift-time structure (board-C mesh defect), so the time-fit plateau is unusable
and only the head-on signature-regression band is quoted. Spark = \% of
crossings coinciding with a discharge.}
\end{keybox}

The right way to read this table:
\begin{itemize}
  \item \textbf{A and B are excellent} and consistent with each other ---
  high uniform efficiency at sub-pitch resolution. Where both were probed
  deeply (charge sharing, charge balance, timing, tracking), they agree,
  which is evidence the performance is a property of the design.
  \item \textbf{C and D reconstruct beautifully when they fire} (note their
  position resolutions --- C's 0.45\,mm is the best core in the fleet) but
  were operated with a spark rate of 23--33\% of crossings during the June
  tests; their HV scans show classic turn-on/plateau/falloff with optima at
  480\,V and 440\,V respectively. Their low headline efficiencies are an
  operating-point effect (drift 700\,V, amplification at/past the sparking
  edge), not broken chambers.
  \item \textbf{E fires on $\sim$70\% of muons but reconstructs only $\sim$21\%}:
  its clusters are too small (too little gas gain) to pass reconstruction.
  This is a gain problem (HV/threshold/gas), the same pathology another
  chamber showed earlier and recoverable in principle --- not a dead detector.
\end{itemize}

\begin{figure}[htbp]
  \centering
  \pkgfig[width=0.85\textwidth]{21-det3A-efficiency-breakdown.png}\\[4pt]
  \pkgfig[width=0.85\textwidth]{25-det7D-efficiency-breakdown.png}
  \caption{The fleet story in two loss budgets: Detector A (top) loses almost
  nothing; Detector D (bottom) fires on the muons just as reliably but loses a
  third of its crossings to discharges at its June operating point.
  Per-detector efficiency maps for all five chambers are in \texttt{figures/}
  (\texttt{2x} series) and the per-detector QA pages in
  \texttt{source\_reports/june-detectors-overview.pdf}.}
  \label{fig:fleetmaps}
\end{figure}

%=====================================================================
\section{Quotable numbers, with the honest caveats}
\label{sec:quotable}

\begin{center}
\small
\begin{tabular}{@{}p{5.2cm}p{3.6cm}p{6.2cm}@{}}
\toprule
Quantity & Value & Fine print \\
\midrule
Active area & $\sim$40$\times$40\,cm$^2$ & 512 strips $\times$ 0.78\,mm, per coordinate \\
Strip pitch & 0.78\,mm & both coordinates \\
Drift gap & 30\,mm & operated at $\approx$330\,V/cm \\
Gas & Ar/iC$_4$H$_{10}$ 95/5 & atmospheric, flushed \\
Drift velocity (nominal) & $34\pm1.5\,\mu$m/ns & measured; matches Magboltz w/ $\sim$1\% H$_2$O \\
Efficiency (A, active area) & $92.9\pm0.2$\,\% & within 5\,mm, operating (spark-free $96.6$\,\%) \\
Efficiency (A, core) & $\sim$96\,\% & $>$25\,mm from frame (pre-2026-07-14, not yet re-verified) \\
Muons producing any signal (A) & 100.0\,\% & chamber essentially never blind (silent $\approx0$\,\%) \\
Position resolution (A) & 0.44--0.47\,mm & sub-pitch; includes telescope error (conservative) \\
Angular resolution (A) & $\approx$1.7$^\circ$ & hybrid tracking, per plane, all angles; bias $\approx0.1^\circ$ \\
Time resolution (A) & 33\,ns ($\approx$1\,mm drift) & telescope-free, inter-plane method \\
Absolute event time (A) & $\sigma_{68}=37.7$\,ns & vs.\ scintillator trigger \\
X/Y charge balance & $0.49$--$0.53$, flat & pixel layer routes charge evenly; $r\approx0.84$--$0.89$ \\
Charge sharing (design) & 45--52\% to neighbour & measured on A and B; enables sub-pitch \\
Spark rate (A, at optimum) & 3.9\,\% of crossings & Poisson 0.33\,Hz; edge-seeded; muon-induced \\
Dead time after spark & none measurable & efficiency \& gain flat vs.\ time since spark \\
Optimal amplification HV & 440--495\,V & per chamber; plateaus $\sim$93\% (A/B) \\
\bottomrule
\end{tabular}
\end{center}

\paragraph{Caveats to keep the talk honest.}
\begin{enumerate}
  \item All results are \textbf{preliminary} (June 2026 bench campaign;
  analysis mature but pre-publication).
  \item Position resolutions are \textbf{detector $\oplus$ telescope}: the
  reference tracker's pointing error is not yet deconvolved, so quoted values
  are upper limits (the planned deconvolution will only improve them).
  \item Results are for \textbf{cosmic muons} (minimum-ionizing, broad angular
  spread) at the specific gas and HV listed; the experiment's particles and
  settings differ in detail.
  \item C, D and E numbers reflect their June \emph{operating points};
  their HV scans indicate better performance is available at optimum settings.
  \item Angular resolutions are quoted \textbf{only from the current-best
  (hybrid) tracking}, run \textbf{fleet-wide} (A 1.7$^\circ$, B 2.3$^\circ$,
  C 3.9$^\circ$, D 3.4$^\circ$, E 2.6$^\circ$; $|\theta|<5^\circ$ band). The 700\,V
  chambers C/D and gain-limited E are degraded but honestly measured. \textbf{C's
  plateau time-fit is unusable} --- its X-plane (mx17\_6, FEU 3) develops no
  micro-TPC drift-time structure (a board-C mesh defect), so only C's low-angle
  signature-regression band is quoted; the plateau is omitted, not averaged in.
\end{enumerate}

%=====================================================================
\appendix
\section{Where each number comes from}

Every value in this document traces to the June-2026 QA analysis
(\texttt{mx\_june\_cosmic\_qa/} in the analysis repository):
\texttt{JUNE\_RESULTS\_SUMMARY.md} (fleet efficiencies, resolutions, HV scans),
\texttt{PAPER\_STATUS.md} and \texttt{MICROTPC\_RUNBOOK.md} (micro-TPC chain:
drift velocity, sharing, angles, position estimators, edge, gas),
the standalone reports in \texttt{source\_reports/} (time resolution, X/Y
charge balance, spark studies), and \texttt{DET3\_WEEKEND\_ANALYSIS.md}
(detector A deep-dive). The reference tracking is the v2 reprocessing (cleaner
track selection), most recently re-tightened 2026-07-14 to $\chi^2<1.0$,
$N_\mathrm{clus}=4$ (was $\chi^2<5$, $N_\mathrm{clus}\ge3$) after finding the
position residual was reference-limited, not chamber-limited --- see
\texttt{det3\_recofar\_analysis/M3\_CUT\_AND\_ACTIVE\_AREA\_NOTE.md}; all five
detectors were re-analyzed with each recipe in turn.

\section{Mini-glossary}
\begin{description}
  \item[Micromegas] Gaseous detector with a micro-mesh separating a drift
  region from a thin amplification gap.
  \item[Resistive strips] A resistive layer over the readout that limits
  discharge currents --- the spark-protection scheme validated in
  Section~\ref{sec:sparks}.
  \item[micro-TPC mode] Using electron drift time as a depth coordinate, so a
  single chamber measures a 3D track segment, not just a 2D point.
  \item[$\sigma_{68}$] Half-width containing 68\% of a distribution ---
  like a Gaussian $\sigma$ but robust to tails.
  \item[Efficiency plateau] HV range where efficiency is maximal and flat;
  the natural operating point.
  \item[Attachment] Capture of drifting electrons by electronegative
  contaminants (O$_2$); shortens the visible charge column.
\end{description}

\end{document}
